\section*{Chapter 2 --- Limits and Continuity}

\begin{solution}{exo:tle:limcont:1}
Dividing by $x^3$:
$\dfrac{2x^3 - x + 1}{5x^3 + x^2} = \dfrac{2 - 1/x^2 + 1/x^3}{5 + 1/x}
\to \dfrac{2}{5}$.

$x^2 + 3x - 1 = x^2(1 + 3/x - 1/x^2) \to +\infty$ as $x \to -\infty$, since
$x^2 \to +\infty$ and the bracket tends to $1$.

As $x \to 2^+$, the numerator tends to $3 > 0$ and the denominator to $0^+$,
so $\dfrac{x+1}{x-2} \to +\infty$.

For $x > 0$, $-\dfrac1x \leq \dfrac{\cos x}{x} \leq \dfrac1x$, so the limit
is $0$ by the squeeze theorem.
\end{solution}

\begin{solution}{exo:tle:limcont:2}
\emph{1.} $2x^2 - 3x + 1 = (x-1)(2x-1)$ (the root $x=1$ is apparent from
$2 - 3 + 1 = 0$). For $x \neq 1$, cancel $x - 1$: $f(x) = 2x - 1$.

\emph{2.} Hence $\lim_{x\to1} f(x) = 1$. Setting $f(1) = 1$ extends $f$ to
the polynomial function $x \mapsto 2x-1$ on $\R$, which is continuous at $1$.
\end{solution}

\begin{solution}{exo:tle:limcont:3}
Write $f(x) = \dfrac{3x - 1}{x + 2} = \dfrac{3(x+2) - 7}{x+2}
= 3 - \dfrac{7}{x+2}$. As $x \to \pm\infty$, $f(x) \to 3$: the line $y = 3$
is a horizontal asymptote at both infinities. As $x \to (-2)^\pm$,
$f(x) \to \mp\infty$: the line $x = -2$ is a vertical asymptote. The curve is
a hyperbola, below $y=3$ for $x > -2$ and above it for $x < -2$.
\end{solution}

\begin{solution}{exo:tle:limcont:4}
Multiply by the conjugate:
\[
\sqrt{x^2 + x} - x = \frac{x^2 + x - x^2}{\sqrt{x^2+x} + x}
= \frac{x}{x\left(\sqrt{1 + 1/x} + 1\right)}
= \frac{1}{\sqrt{1 + 1/x} + 1}
\xrightarrow[x\to+\infty]{} \frac{1}{2}.
\]
Similarly,
\[
\frac{\sqrt{1+x} - 1}{x} = \frac{(1+x) - 1}{x\left(\sqrt{1+x} + 1\right)}
= \frac{1}{\sqrt{1+x} + 1} \xrightarrow[x\to0]{} \frac{1}{2}.
\]
\end{solution}

\begin{solution}{exo:tle:limcont:5}
\emph{1.} Since $\sin x \geq -1$, $f(x) \geq x - 1 \to +\infty$, and we
conclude by comparison. Yet $f'(x) = 1 + \cos x$ vanishes at every odd
multiple of $\pi$, and on every interval $\intco{A}{+\infty}$ the function
alternates between intervals of strict increase; it is increasing but not
strictly, and more importantly its growth is not needed: comparison suffices.

\emph{2.} For $x > 1$, divide by $x$:
\[
g(x) = \frac{1 + \frac{\sin x}{x}}{1 - \frac{\sin x}{x}},
\]
and $\frac{\sin x}{x} \to 0$ by the squeeze theorem
($\abs{\sin x} \leq 1$). Hence $g(x) \to \frac{1+0}{1-0} = 1$.
\end{solution}

\begin{solution}{exo:tle:limcont:6}
$f(x) = x^3 - 3x + 1$ is continuous, with $f'(x) = 3x^2 - 3 = 3(x-1)(x+1)$:
$f$ increases on $\intoc{-\infty}{-1}$, decreases on $\intcc{-1}{1}$,
increases on $\intco{1}{+\infty}$, with local maximum $f(-1) = 3 > 0$ and
local minimum $f(1) = -1 < 0$, $\lim_{-\infty} f = -\infty$,
$\lim_{+\infty} f = +\infty$.

On each of the three monotonicity intervals, $f$ is continuous, strictly
monotonic, and $0$ lies strictly between the endpoint values/limits, so the
bijection theorem gives exactly one zero per interval: three zeros in total.
Sign checks: $f(-2) = -1 < 0 < f(-1) = 3$, so one zero in
$\intoo{-2}{-1}$; $f(0) = 1 > 0 > f(1) = -1$, one zero in $\intoo{0}{1}$;
$f(1) = -1 < 0 < f(2) = 3$, one zero in $\intoo{1}{2}$.
\end{solution}

\begin{solution}{exo:tle:limcont:7}
Let $g(x) = f(x) - x$, continuous on $\intcc{0}{1}$. Since
$f(0) \in \intcc{0}{1}$, $g(0) = f(0) \geq 0$; since $f(1) \in \intcc{0}{1}$,
$g(1) = f(1) - 1 \leq 0$. By the intermediate value theorem, $g$ vanishes at
some $c \in \intcc{0}{1}$, i.e.\ $f(c) = c$.
\end{solution}

\begin{solution}{exo:tle:limcont:8}
Let $u(t)$ and $d(t)$ be the positions (measured as distance from the bottom
of the trail) at time $t \in \intcc{8}{20}$ on the way up and on the way
down. Both are continuous, $u(8) = 0$, $u(20) = L$ (the length of the trail),
$d(8) = L$, $d(20) = 0$. The function $g = u - d$ is continuous with
$g(8) = -L \leq 0$ and $g(20) = L \geq 0$, so $g(c) = 0$ for some time $c$ by
the intermediate value theorem: at time $c$, both days' positions coincide.
\end{solution}

\begin{solution}{exo:tle:limcont:9}
\emph{1.} $f'(x) = 5x^4 + 3x^2 + 1 > 0$, so $f$ is strictly increasing on
$\R$; with $\lim_{\pm\infty} f = \pm\infty$ and continuity, the bijection
theorem gives a unique real zero $\alpha$. Since $f(0) = -1 < 0$ and
$f(1) = 2 > 0$, $\alpha \in \intoo{0}{1}$.

\emph{2.} Maintain an interval $\intcc{a_n}{b_n}$ with
$f(a_n) < 0 < f(b_n)$, starting from $\intcc{0}{1}$; at each step evaluate
$f$ at the midpoint $m$ and keep the half-interval on which the sign change
persists. After $n$ steps the interval has length $2^{-n}$, so the midpoint
approximates $\alpha$ to within $2^{-(n+1)}$. We need
$2^{-(n+1)} \leq 10^{-6}$, i.e.\ $n + 1 \geq \frac{6}{\log_{10} 2} \approx 19.93$:
$n = 20$ iterations suffice.
\end{solution}
